\section{The requested theorem and the actual result}

\subsection{The mathematical statement}

Fix a number field $K$, its ring of integers $\mathcal O_K$, and a smooth
geometrically integral affine curve $C$ over $K$ with a specified embedding
in affine space. An integral point is a tuple in $\mathcal O_K^n$ satisfying
the defining equations of that affine presentation. Let $\overline C$ be
the smooth projective completion and let $g$ be its geometric genus.
Siegel's theorem states that
\[
 g\ge1\quad\Longrightarrow\quad C(\mathcal O_K)\text{ is finite}.
\]
The more general Siegel--Mahler statement permits a finite set $S$ of places
containing the infinite places and the ring $\mathcal O_{K,S}$ of elements
integral outside $S$. If $r$ is the number of geometric points in
$\overline C\setminus C$, the sufficient condition is $g\ge1$ or $r\ge3$.
An affine curve has $r\ge1$, so this condition is equivalent to
$2g-2+r>0$. We retain the ordinary-integral, positive-genus statement as
the requested goal. The classical background is discussed in
\cite{BiluIllengo2009,CorvajaZannier2002}.

The embedding and the deleted points matter. A birational equivalence of
smooth projective curves can introduce denominators, and it generally does
not take the original integral-point problem to the integral-point problem
of an arbitrary new affine equation. Thus replacing a curve by a convenient
equation requires a proved, effective comparison of the integral models.

For example, the existence of hyperelliptic equations for genus-two curves
does not settle all their integral-point problems. Corvaja--Lombardo--Zannier
explain precisely this distinction for curves with one point removed, and
report that a general terminating procedure for their integral-point sets
is not known \cite[Introduction and Remark 1.2]{CorvajaLombardoZannier2024}.
Our explicit results concern their stated equations over $\Z$. They do not
erase this dependence on the affine presentation.

\subsection{Why a full affirmative answer is not claimed}

The usual classical theorem does not, in general, provide a computable bound
on the sizes of the points. A constructive theorem producing a finite
complete list, or a bound followed by finite search, would supply that
missing information. In iPRA the witness consequence is particularly
strong: under decidable input hypotheses a uniform proof would yield a
primitive recursive bound. Section~\ref{sec:logical-audit} proves this
conditional implication carefully.

This is a serious obstacle, not a proof of impossibility. We neither prove
that Siegel's theorem is independent of iPRA nor claim that the effective
problem cannot be solved. Moreover, the user did not specify that
``finite'' must mean a positive complete list. A bound on the \emph{number}
of points is a different meaningful constructive target. That target does
not automatically settle the height problem. A complete iPRA proof of the
general cardinality formulation is also absent here.

The recent primary source above and the older effective results describe
many cases with computable solutions. This report does not claim that
positive genus itself forces ineffectivity, or that all genus-two models
are beyond effective methods. The open general step and the elementary
families below must be distinguished.

\subsection{Results established in the report}

The main arithmetic theorem applies to a checkable finite certificate,
not to a promised numerical bound. Write $A(X)=a_mX^m+\cdots+a_0$ with
$a_m\ne0$, and let $R(X)=\sum_i r_iX^i$. Suppose
\[
 D^2f(X)=A(X)^2+R(X),\qquad
 D\in\Z_{>0},\quad m\ge1,\quad\deg R<m,\quad R\ne0.
\]
Then every integral point $y^2=f(x)$ satisfies
\[
 |x|\le B_x:=\sum_{i<m}|a_i|+\sum_i|r_i|,
 \qquad
 |y|\le B_y:=\sum_i |f_i|\max(1,B_x)^i.
\]
These are explicit integer bounds, and both the certificate and the
enumeration can be computed by primitive recursion. The hypotheses are
verified coefficient by coefficient. In particular, the certificate
does not contain an assertion about all solutions as an unchecked field.

We give a uniform certificate constructor for every nonsquare integral
polynomial $f$ of even positive degree whose leading coefficient is a
nonzero integer square. When $f$ is squarefree of degree $2m\ge4$, this
gives a genuine positive-genus instance of Siegel finiteness. The proof
is an elementary instance of Runge's method: two integral quantities
cannot be nonzero factors of a smaller integer.

We additionally give a sharper factor-divisor algorithm for
$y^2=x^4+ax^2+b$, a complete six-point classification for $a=7,b=1$,
and a four-point classification valid in every positive genus. No
claim of novelty is made for the underlying factorization or Runge
principle; the purpose is an explicit arithmetic proof and a transparent
formalization interface.

\subsection{Repository snapshot and reuse}

The repository was inspected through the GitHub connector at revision
\begin{center}\small\ttfamily
7eeb3c0e821594d5d5ebcfbaf4bdba64819e3502.
\end{center}
The actual foundation supplies primitive recursive syntax, intuitionistic
natural deduction, quantifier-free induction, and faithful translation
to the concise native calculus. We use those interfaces in
Section~\ref{sec:lean}; we do not infer an arithmetic theorem from
bi-interpretability alone.

The repository's proof-packet index explicitly says that the packets were
filed as delivered, without being rerun or verified at placement.
The inspected Faltings package explicitly leaves general Faltings in iPRA
unproved; the Mordell--Weil package retains a covering-list hypothesis in
its general descent theorem; and the Roth package leaves a general
effective cutoff unproved. Their existence therefore supplies no verified
general Siegel height bound. These observations are limited to the
identified files and are not an exhaustive correctness review of ProveIt.
Exact paths and the pinned source manifest accompany the archive
\cite{ProveIt2026}.

\subsection{Status at a glance}

\begin{center}
\begin{tabular}{@{}p{0.48\linewidth}p{0.43\linewidth}@{}}
\toprule
Claim & Status in this archive\\
\midrule
General positive Siegel finiteness in iPRA & Not proved\\
General cardinality Siegel finiteness in iPRA & Not proved\\
General negative translation in iPRA & Not proved\\
Explicit Runge bound and certificate constructor & Fully developed arithmetic proof\\
Exhaustive search for the certified family & Proved; executable reference code\\
Every-genus four-point family & Complete arithmetic classification\\
Biquadratic factor bound and six-point example & Written proof and compiled native Lean proof\\
Embedded \texttt{Prf PRAAx} derivations & Detailed implementation plan; not constructed\\
\bottomrule
\end{tabular}
\end{center}
